% when: pgf, pgfcore, pgfmath, tikz, pgfplots, beamer, tcolorbox, tikz-cd, circuitikz, forest, todonotes, pgfgantt, smartdiagram, chemfig, tikzposter, pgf-pie, tkz-euclide, bchart, tikzpagenodes
% PGF evaluates every length it is given (\pgfpoint{3bp}{-1bp}) with pgfmath's expression parser, which takes hundreds
% of steps even for a plain length. Beamer and TikZ pictures give it the same few plain ones thousands of times: a
% length like 8bp or -1.5pt, or a length register. Those get the result pgfmath's parser gives, computed the way the
% parser computes it (the same TeX assignment, the same units flags and \pgfmathpostparse). An expression's kind is
% found from its text the first time pgfmath's own number check passes it to the parser, and kept for the first 2000
% such expressions. In Clearsnap's engine, every other expression the parser is handed goes to the engine's own
% evaluator first (engine/native/pgfmath.c), which gives pgfmath's result for plain arithmetic. Anything else, and
% everything while the fpu library is on, runs pgfmath's own code. Installed at \begin{document}, over the pgfmath it
% was written against only.
\makeatletter
\def\clearsnap@pm@expect{%
  \begingroup
    \pgfmath@catcodes
    \pgfmath@quickparsefalse
    \ifpgfmathfloat
      \let\pgfmath@parse@next=\pgfmathfloatparsenumber
    \else
      \let\pgfmath@parse@next=\pgfmathparse@
    \fi
    \pgfmath@parse@next}
\def\clearsnap@pm@expect@#1{%
    \global\pgfmathunitsdeclaredfalse
    \global\pgfmathmathunitsdeclaredfalse
    \pgfmath@protected@edef\pgfmath@expression{#1}%
    \expandafter\pgfmathparse@trynumber@loop\pgfmath@expression\pgfmath@parse@stop
}
\def\clearsnap@pm@expectnum{%
    \let\pgfmathresult=\pgfmath@expression
    \pgfmath@smuggleone\pgfmathresult%
  \endgroup
  \ignorespaces
}
\def\clearsnap@pm@expectend#1\pgfmath@parse@stop{%
    \pgfmathpostparse
    \pgfmath@smuggleone\pgfmathresult%
  \endgroup
  \ignorespaces}
\def\clearsnap@pm@expectneg#1{%
    \begingroup%
        \pgfmath@x=-#1pt\relax%
        \edef\pgfmathresult{\pgfmath@tonumber{\pgfmath@x}}%
        \pgfmath@smuggleone\pgfmathresult%
    \endgroup%
}
\def\clearsnap@pm@expectadd#1#2{%
    \begingroup%
    \pgfmath@x=#1pt\relax%
    \pgfmath@y=#2pt\relax%
    \advance\pgfmath@x by\pgfmath@y%
    \pgfmath@returnone\pgfmath@x%
  \endgroup%
}
\def\clearsnap@pm@expectsubtract#1#2{%
    \begingroup%
    \pgfmath@x=#1pt\relax%
    \pgfmath@y=#2pt\relax%
    \advance\pgfmath@x by-\pgfmath@y%
    \pgfmath@returnone\pgfmath@x%
  \endgroup%
}
% pgf 3.1.10 and 3.1.11a (whose parser only adds \string and \noexpand around tokens it looks up) parse lengths alike.
\def\clearsnap@pm@versions{,3.1.10,3.1.11a,}
\def\clearsnap@pm@install{%
  \ifx\pgfmathparse\clearsnap@pm@expect \ifx\pgfmathparse@\clearsnap@pm@expect@
  \ifx\pgfmathparse@expression@is@number\clearsnap@pm@expectnum \ifx\pgfmath@parse@end\clearsnap@pm@expectend
    \edef\clearsnap@pm@tmp{\noexpand\in@{,\pgfversion,}{\clearsnap@pm@versions}}\clearsnap@pm@tmp
    \edef\clearsnap@pm@tmp{\csname pgfmath@operation@+@name\endcsname/\csname pgfmath@operation@-@name\endcsname/%
      \csname pgfmath@operation@+@precedence\endcsname/\csname pgfmath@operation@-@precedence\endcsname/%
      \csname pgfmath@operation@neg@precedence\endcsname}%
    \def\clearsnap@pm@ops{add/subtract/500/500/500}%
    \ifx\clearsnap@pm@tmp\clearsnap@pm@ops\else\in@false\fi
    \ifin@\global\let\pgfmathparse\clearsnap@pm@parse\clearsnap@pm@probe\fi
  \fi\fi\fi\fi}
% Clearsnap's engine evaluates arithmetic and the common functions itself (engine/native/pgfmath.c): \pdfstrcmp
% (XeTeX: \strcmp) with this first argument and pgfmath's trig format gives U or N (units or not) and pgfmath's
% result, or F to leave the expression to pgfmath.
% Other engines answer the probe with a comparison, and the native step is skipped.
\def\clearsnap@pm@probe{%
  \global\let\clearsnap@pm@native\@empty
  \edef\clearsnap@pm@ok{\detokenize{N2.0}}%
  \ifdefined\pdfstrcmp\global\let\clearsnap@pm@cmp\pdfstrcmp\clearsnap@pm@probe@\fi
  \ifx\clearsnap@pm@native\@empty\ifdefined\strcmp\global\let\clearsnap@pm@cmp\strcmp\clearsnap@pm@probe@\fi\fi
  \ifx\clearsnap@pm@native\relax\clearsnap@fpu@install\fi}
\def\clearsnap@pm@probe@{%
  \edef\clearsnap@pm@tmp{\clearsnap@pm@cmp{clearsnapnativepgfmath0}{1+1}}%
  \ifx\clearsnap@pm@tmp\clearsnap@pm@ok\global\let\clearsnap@pm@native\relax\fi}
\let\clearsnap@pm@native\@empty
\def\clearsnap@pm@natsplit#1#2\relax{\def\clearsnap@pm@natkind{#1}\def\clearsnap@pm@natval{#2}}
\def\clearsnap@pm@natdone{%
  \if U\clearsnap@pm@natkind\global\pgfmathunitsdeclaredtrue\fi
  \let\pgfmathresult\clearsnap@pm@natval
  \pgfmath@parse@end\pgfmath@parse@stop}
% pgfmath's fpu library (pgfplots computes with it): its parse step and the float operations pgfplots calls most go to
% the engine's evaluator first, and F leaves them to the library's own code (\pgfmathfloatparsenumber, add,
% subtract, multiply, divide, tofixed, neg, abs, ln, sqrt). pgfplots' x and y (zero-argument functions
% that give \x and \y) are passed along when they are exactly those. Installed with the probe above, over the fpu
% library of the same pgf versions only, while its extended precision is the default 2.
\edef\clearsnap@pm@F{\detokenize{F}}
\def\clearsnap@fpu@expect#1#2\pgfmathfloat@{%
    \def\pgfmathfloat@test{#1}%
    \ifx\pgfmathfloat@test\pgfmathfloat@char@asterisk%
        \def\pgfmathfloat@expression{#2}%
        \pgfmathfloat@scaleactivetrue
    \fi%
    \expandafter\pgfmathparse@\expandafter{\pgfmathfloat@expression}%
}%
% The engine's floats are other-character text; a float's Y and e are letters.
\begingroup\lccode`\!=`\Y \lccode`\?=`\e
\lowercase{\endgroup\def\clearsnap@fpu@mk#1!#2?#3]#4\relax}{\def\pgfmathresult{#1Y#2e#3]}}
\def\clearsnap@fpu@install{%
  \edef\clearsnap@pm@tmp{\clearsnap@pm@cmp{clearsnapnativefpuop:parse}{2}}%
  \edef\clearsnap@pm@ok{\detokenize{1Y2.0e0]}}%
  \ifx\clearsnap@pm@tmp\clearsnap@pm@ok \ifx\pgfmathfloatparse@@\clearsnap@fpu@expect
    \global\let\pgfmathfloatparse@@\clearsnap@fpu@parse@@
    \clearsnap@fpu@put\pgfmathfloatparsenumber\pgfmathfloatqparsenumber{parse}%
    \clearsnap@fpu@put\pgfmathfloattofixed@\pgfmathfloattofixed{tofixed}%
    \clearsnap@fpu@put\pgfmathfloatneg@\relax{neg}%
    \clearsnap@fpu@put\pgfmathfloatabs@\relax{abs}%
    \clearsnap@fpu@put\pgfmathfloatln@\pgfmathfloatln{ln}%
    \clearsnap@fpu@put\pgfmathfloatsqrt@\pgfmathfloatsqrt{sqrt}%
    \clearsnap@fpu@put\pgfmathfloatadd@\pgfmathfloatadd{add}%
    \clearsnap@fpu@put\pgfmathfloatsubtract@\pgfmathfloatsubtract{sub}%
    \clearsnap@fpu@put\pgfmathfloatmultiply@\pgfmathfloatmultiply{mul}%
    \clearsnap@fpu@put\pgfmathfloatdivide@\pgfmathfloatdivide{div}%
  \fi\fi}
% #1 the library's macro, #2 its other name (if still the same macro), #3 the operation: the original is kept as
% \clearsnap@fpu@o@<operation> and #1 (and #2) become \clearsnap@fpu@<operation>.
\def\clearsnap@fpu@put#1#2#3{%
  \expandafter\global\expandafter\let\csname clearsnap@fpu@o@#3\endcsname#1%
  \ifx#2#1\expandafter\global\expandafter\let\expandafter#2\csname clearsnap@fpu@#3\endcsname\fi
  \expandafter\global\expandafter\let\expandafter#1\csname clearsnap@fpu@#3\endcsname}
\def\clearsnap@fpu@one#1#2{\edef\clearsnap@pm@natres{\clearsnap@pm@cmp{clearsnapnativefpuop:#1}{#2}}}
\def\clearsnap@fpu@float{\expandafter\clearsnap@fpu@mk\clearsnap@pm@natres\relax}
\def\clearsnap@fpu@float@#1{\clearsnap@fpu@float}
\def\clearsnap@fpu@float@@#1#2{\clearsnap@fpu@float}
\def\clearsnap@fpu@text@#1{\let\pgfmathresult\clearsnap@pm@natres}
\def\clearsnap@fpu@parse#1{\clearsnap@fpu@one{parse}{#1}%
  \ifx\clearsnap@pm@natres\clearsnap@pm@F\expandafter\clearsnap@fpu@o@parse\else\expandafter\clearsnap@fpu@float@\fi{#1}}
\def\clearsnap@fpu@tofixed#1{\clearsnap@fpu@one{tofixed}{#1}%
  \ifx\clearsnap@pm@natres\clearsnap@pm@F\expandafter\clearsnap@fpu@o@tofixed\else\expandafter\clearsnap@fpu@text@\fi{#1}}
\def\clearsnap@fpu@neg#1{\clearsnap@fpu@one{neg}{#1}%
  \ifx\clearsnap@pm@natres\clearsnap@pm@F\expandafter\clearsnap@fpu@o@neg\else\expandafter\clearsnap@fpu@float@\fi{#1}}
\def\clearsnap@fpu@abs#1{\clearsnap@fpu@one{abs}{#1}%
  \ifx\clearsnap@pm@natres\clearsnap@pm@F\expandafter\clearsnap@fpu@o@abs\else\expandafter\clearsnap@fpu@float@\fi{#1}}
\def\clearsnap@fpu@ln#1{\clearsnap@fpu@one{ln}{#1}%
  \ifx\clearsnap@pm@natres\clearsnap@pm@F\expandafter\clearsnap@fpu@o@ln\else\expandafter\clearsnap@fpu@float@\fi{#1}}
\def\clearsnap@fpu@sqrt#1{\clearsnap@fpu@one{sqrt}{#1}%
  \ifx\clearsnap@pm@natres\clearsnap@pm@F\expandafter\clearsnap@fpu@o@sqrt\else\expandafter\clearsnap@fpu@float@\fi{#1}}
\def\clearsnap@fpu@add#1#2{%
  \ifnum\pgfmathfloatextprec@shift=2 \clearsnap@fpu@one{add}{#1,#2}\else\let\clearsnap@pm@natres\clearsnap@pm@F\fi
  \ifx\clearsnap@pm@natres\clearsnap@pm@F\expandafter\clearsnap@fpu@o@add\else\expandafter\clearsnap@fpu@float@@\fi{#1}{#2}}
\def\clearsnap@fpu@sub#1#2{%
  \ifnum\pgfmathfloatextprec@shift=2 \clearsnap@fpu@one{sub}{#1,#2}\else\let\clearsnap@pm@natres\clearsnap@pm@F\fi
  \ifx\clearsnap@pm@natres\clearsnap@pm@F\expandafter\clearsnap@fpu@o@sub\else\expandafter\clearsnap@fpu@float@@\fi{#1}{#2}}
\def\clearsnap@fpu@mul#1#2{\clearsnap@fpu@one{mul}{#1,#2}%
  \ifx\clearsnap@pm@natres\clearsnap@pm@F\expandafter\clearsnap@fpu@o@mul\else\expandafter\clearsnap@fpu@float@@\fi{#1}{#2}}
\def\clearsnap@fpu@div#1#2{\clearsnap@fpu@one{div}{#1,#2}%
  \ifx\clearsnap@pm@natres\clearsnap@pm@F\expandafter\clearsnap@fpu@o@div\else\expandafter\clearsnap@fpu@float@@\fi{#1}{#2}}
% The parse step: pgfmath's \pgfmathparse@ as the library calls it, the engine first.
\def\clearsnap@fpu@parse@@#1#2\pgfmathfloat@{%
  \def\pgfmathfloat@test{#1}%
  \ifx\pgfmathfloat@test\pgfmathfloat@char@asterisk
    \def\pgfmathfloat@expression{#2}%
    \pgfmathfloat@scaleactivetrue
  \fi
  \expandafter\clearsnap@fpu@parse@\expandafter{\pgfmathfloat@expression}}
\def\clearsnap@fpu@xexpect{\edef\pgfmathresult{\tikz@plot@var}}
\def\clearsnap@fpu@yexpect{\edef\pgfmathresult{\pgfplots@plot@var@y}}
% (pgfplots binds x before \x exists, and pgfmath only reads \x when the expression has an x, so a binding is the text
% \x stands for, when it is defined; anything that is not a float or number there leaves the expression to pgfmath.)
\def\clearsnap@fpu@binds{%
  \ifx\pgfmathx@\clearsnap@fpu@xexpect\expandafter\clearsnap@fpu@bind\tikz@plot@var x\fi
  \ifx\pgfmathy@\clearsnap@fpu@yexpect\expandafter\clearsnap@fpu@bind\pgfplots@plot@var@y y\fi}
\def\clearsnap@fpu@bind#1#2{\ifdefined#1\ifx#1\relax\else;#2=\detokenize\expandafter{#1}\fi\fi}
\def\clearsnap@fpu@parse@#1{%
  \global\pgfmathunitsdeclaredfalse
  \global\pgfmathmathunitsdeclaredfalse
  \pgfmath@protected@edef\pgfmath@expression{#1}%
  \let\clearsnap@pm@kind\clearsnap@pm@slow
  \ifnum\pgfmathfloatextprec@shift=2
    \edef\clearsnap@pm@natres{\clearsnap@pm@cmp{clearsnapnativefpu\pgfmath@trig@format@choice\clearsnap@fpu@binds}%
      {\unexpanded\expandafter{\pgfmath@expression}}}%
    \expandafter\clearsnap@pm@natsplit\clearsnap@pm@natres\relax
    \if F\clearsnap@pm@natkind\else\let\clearsnap@pm@kind\clearsnap@fpu@done\fi
  \fi
  \clearsnap@pm@kind}
\def\clearsnap@fpu@done{%
  \if U\clearsnap@pm@natkind\global\pgfmathunitsdeclaredtrue\fi
  \expandafter\clearsnap@fpu@mk\clearsnap@pm@natval\relax
  \pgfmath@parse@end\pgfmath@parse@stop}
% pgfmath's \pgfmathparse with its own \pgfmathparse@ in place of pgfmath's.
\def\clearsnap@pm@parse{%
  \begingroup
    \pgfmath@catcodes
    \pgfmath@quickparsefalse
    \ifpgfmathfloat
      \let\pgfmath@parse@next=\pgfmathfloatparsenumber
    \else
      \let\pgfmath@parse@next=\clearsnap@pm@parse@
    \fi
    \pgfmath@parse@next}
% pgfmath's \pgfmathparse@: a kept kind runs at once; otherwise pgfmath's number check runs, and if it hands the
% expression to the parser (\pgfmathparse@@, which parses \pgfmath@expression), \clearsnap@pm@full sees it first.
\def\clearsnap@pm@parse@#1{%
  \global\pgfmathunitsdeclaredfalse
  \global\pgfmathmathunitsdeclaredfalse
  \pgfmath@protected@edef\pgfmath@expression{#1}%
  \let\clearsnap@pm@kind\clearsnap@pm@slow
  \ifnum\escapechar=92
    \edef\clearsnap@pm@key{\detokenize\expandafter{\pgfmath@expression}}%
    \ifcsname clearsnap@pm@k@\clearsnap@pm@key\endcsname
      \expandafter\let\expandafter\clearsnap@pm@kind\csname clearsnap@pm@k@\clearsnap@pm@key\endcsname
    \else
      \let\clearsnap@pm@real\pgfmathparse@@
      \let\pgfmathparse@@\clearsnap@pm@full
    \fi
  \fi
  \clearsnap@pm@kind}
\def\clearsnap@pm@stop{\clearsnap@pm@stop}
\def\clearsnap@pm@slow{\expandafter\pgfmathparse@trynumber@loop\pgfmath@expression\pgfmath@parse@stop}
\def\clearsnap@pm@parser{\pgfmathparse@@\pgfmath@parse@stop}
\newcount\clearsnap@pm@count
\def\clearsnap@pm@full#1\pgfmath@parse@stop{%
  \let\pgfmathparse@@\clearsnap@pm@real
  \ifx\clearsnap@pm@native\relax
    \edef\clearsnap@pm@natres{\clearsnap@pm@cmp{clearsnapnativepgfmath\pgfmath@trig@format@choice}%
      {\unexpanded\expandafter{\pgfmath@expression}}}%
    \expandafter\clearsnap@pm@natsplit\clearsnap@pm@natres\relax
    \if F\clearsnap@pm@natkind\let\clearsnap@pm@kind\clearsnap@pm@parser\else\let\clearsnap@pm@kind\clearsnap@pm@natdone\fi
  \else\ifnum\clearsnap@pm@count<2000
    \global\advance\clearsnap@pm@count\@ne
    \clearsnap@pm@classify
    \ifx\clearsnap@pm@kind\clearsnap@pm@slow\let\clearsnap@pm@kind\clearsnap@pm@parser\fi
  \else
    \let\clearsnap@pm@kind\clearsnap@pm@parser
  \fi\fi
  \clearsnap@pm@kind}
% A length (8bp, -1.5pt, \pgf@x, -\pgf@x): the parser assigns it to \pgfmath@dimen as TeX reads it, and the result is
% that length as a number (with a minus, pgfmath's neg of it: the same number negated).
\def\clearsnap@pm@length{%
  \pgfmath@dimen=\pgfmath@expression\relax
  \global\pgfmathunitsdeclaredtrue
  \edef\pgfmathresult{\pgfmath@tonumber{\pgfmath@dimen}}%
  \pgfmath@parse@end\pgfmath@parse@stop}
\def\clearsnap@pm@split#1#2\clearsnap@pm@stop{\def\clearsnap@pm@first{#1}\def\clearsnap@pm@rest{#2}}
\def\clearsnap@pm@neg{%
  \expandafter\clearsnap@pm@split\pgfmath@expression\clearsnap@pm@stop
  \ifx\clearsnap@pm@first\pgfmath@char@minus \ifx\pgfmathneg@\clearsnap@pm@expectneg
    \let\clearsnap@pm@next\clearsnap@pm@length
  \else\let\clearsnap@pm@next\clearsnap@pm@slow\fi
  \else\let\clearsnap@pm@next\clearsnap@pm@slow\fi
  \clearsnap@pm@next}
% A register is used while it is the dimen register it was when its expression was first seen.
\def\clearsnap@pm@reg{%
  \expandafter\ifx\csname clearsnap@pm@r@\clearsnap@pm@key\expandafter\endcsname\pgfmath@expression
    \expandafter\clearsnap@pm@length
  \else
    \expandafter\clearsnap@pm@slow
  \fi}
\def\clearsnap@pm@negreg{%
  \expandafter\clearsnap@pm@split\pgfmath@expression\clearsnap@pm@stop
  \expandafter\ifx\csname clearsnap@pm@r@\clearsnap@pm@key\expandafter\endcsname\clearsnap@pm@rest
    \expandafter\clearsnap@pm@neg
  \else
    \expandafter\clearsnap@pm@slow
  \fi}
% The kind of an expression the parser was given, kept for the next time: a length, a register (while it is the dimen
% register it was) or "slow" (\clearsnap@pm@slow: pgfmath's number check and parser).
\def\clearsnap@pm@classify{%
  \let\clearsnap@pm@kind\clearsnap@pm@slow
  \edef\clearsnap@pm@tmp{\noexpand\in@{ }{\clearsnap@pm@key}}\clearsnap@pm@tmp
  \ifin@
    \clearsnap@pm@regscan
  \else
    \let\clearsnap@pm@sign\@empty
    \expandafter\clearsnap@pm@sA\clearsnap@pm@key\clearsnap@pm@stop\clearsnap@pm@stop\clearsnap@pm@end
  \fi
  \ifx\clearsnap@pm@kind\clearsnap@pm@slow\clearsnap@pm@chainscan\fi
  \global\expandafter\let\csname clearsnap@pm@k@\clearsnap@pm@key\endcsname\clearsnap@pm@kind}
% A length is an optional -, digits with at most one period (at least one digit), and one of pgfmath's two-letter
% units that do not depend on the font: the text read one character at a time.
\def\clearsnap@pm@end{\clearsnap@pm@end}
\def\clearsnap@pm@pad{\clearsnap@pm@stop\clearsnap@pm@stop}
\def\clearsnap@pm@sX#1\clearsnap@pm@end{}
\def\clearsnap@pm@class#1#2#3#4{% next: after a digit #2, after a period #3, anything else #4 (which reads it again)
  \def\clearsnap@pm@nx{#4#1}%
  \ifx#1\clearsnap@pm@stop\let\clearsnap@pm@nx\clearsnap@pm@sX\else
    \ifnum`#1>47 \ifnum`#1<58 \let\clearsnap@pm@nx#2\fi\fi
    \if.#1\let\clearsnap@pm@nx#3\fi
  \fi}
\def\clearsnap@pm@sA#1{%
  \ifx#1\clearsnap@pm@stop\let\clearsnap@pm@nx\clearsnap@pm@sX\else
    \if-#1\def\clearsnap@pm@sign{-}\let\clearsnap@pm@nx\clearsnap@pm@sB
    \else\clearsnap@pm@class#1\clearsnap@pm@sC\clearsnap@pm@sD\clearsnap@pm@sX\fi
  \fi
  \clearsnap@pm@nx}
\def\clearsnap@pm@sB#1{\clearsnap@pm@class#1\clearsnap@pm@sC\clearsnap@pm@sD\clearsnap@pm@sX\clearsnap@pm@nx}
\def\clearsnap@pm@sC#1{\clearsnap@pm@class#1\clearsnap@pm@sC\clearsnap@pm@sE\clearsnap@pm@sU\clearsnap@pm@nx}
\def\clearsnap@pm@sD#1{\clearsnap@pm@class#1\clearsnap@pm@sE\clearsnap@pm@sX\clearsnap@pm@sX\clearsnap@pm@nx}
\def\clearsnap@pm@sE#1{\clearsnap@pm@class#1\clearsnap@pm@sE\clearsnap@pm@sU\clearsnap@pm@sU\clearsnap@pm@nx}
\def\clearsnap@pm@sU#1#2#3\clearsnap@pm@end{%
  \def\clearsnap@pm@tmp{#3}%
  \ifx\clearsnap@pm@tmp\clearsnap@pm@pad
    \ifcsname clearsnap@pm@u@\string#1\string#2\endcsname
      \ifx\clearsnap@pm@sign\@empty\let\clearsnap@pm@kind\clearsnap@pm@length\else\let\clearsnap@pm@kind\clearsnap@pm@neg\fi
    \fi
  \fi}
\@for\clearsnap@pm@tmp:={pt,bp,cm,mm,in,pc,dd,cc,sp,nd,nc}\do{\expandafter\let\csname clearsnap@pm@u@\clearsnap@pm@tmp\endcsname\@empty}
% A sum or difference of terms (a length, a register with or without a factor, a number), after a minus or not, in
% braces or not (-1, \bmb@boxwidth+4bp, 0.5\bmb@boxwidth+6bp, {64.80573pt}). pgfmath reads each term as TeX reads it,
% negates and adds them up with \advance (its neg, add and subtract), so the code kept for the expression makes the same
% assignments. It is used for this exact token list (\ifx) while pgfmath's neg, add and subtract and the registers are
% what they were.
\def\clearsnap@pm@chain{%
  \let\clearsnap@pm@next\clearsnap@pm@slow
  \expandafter\ifx\csname clearsnap@pm@e@\clearsnap@pm@key\endcsname\pgfmath@expression
    \let\clearsnap@pm@next\clearsnap@pm@chain@
    \csname clearsnap@pm@g@\clearsnap@pm@key\endcsname
  \fi
  \clearsnap@pm@next}
\def\clearsnap@pm@chain@{\csname clearsnap@pm@x@\clearsnap@pm@key\endcsname\pgfmath@parse@end\pgfmath@parse@stop}
\def\clearsnap@pm@chk#1#2{\ifx#1#2\else\let\clearsnap@pm@next\clearsnap@pm@slow\fi}
% The scan reads the expression (pgfmath's braces stripped) a token at a time; a brace or a space ends it.
\def\clearsnap@pm@chainscan{%
  \expandafter\pgfmath@bgroup@strip\expandafter{\pgfmath@expression}%
  \let\clearsnap@pm@code\@empty \let\clearsnap@pm@guard\@empty \let\clearsnap@pm@term\@empty
  \let\clearsnap@pm@sign\@empty \let\clearsnap@pm@op\@empty \let\clearsnap@pm@unit\@empty
  \let\clearsnap@pm@terms\@empty \let\clearsnap@pm@good\relax \def\clearsnap@pm@st{A}%
  \expandafter\clearsnap@pm@read\pgfmath@bgroup@stripped\clearsnap@pm@end
  \ifx\clearsnap@pm@good\relax
    % One unsigned number is pgfmath's number as typed, not a computed one.
    \edef\clearsnap@pm@tmp{\clearsnap@pm@sign\clearsnap@pm@terms}%
    \def\clearsnap@pm@single{n}%
    \ifx\clearsnap@pm@tmp\clearsnap@pm@single\else
      \global\expandafter\let\csname clearsnap@pm@e@\clearsnap@pm@key\endcsname\pgfmath@expression
      \global\expandafter\let\csname clearsnap@pm@g@\clearsnap@pm@key\endcsname\clearsnap@pm@guard
      \expandafter\in@\expandafter u\expandafter{\clearsnap@pm@terms}%
      \expandafter\xdef\csname clearsnap@pm@x@\clearsnap@pm@key\endcsname{\unexpanded\expandafter{\clearsnap@pm@code}%
        \ifin@\global\noexpand\pgfmathunitsdeclaredtrue\fi
        \edef\noexpand\pgfmathresult{\noexpand\pgfmath@tonumber{\noexpand\pgfmath@dimen}}}%
      \let\clearsnap@pm@kind\clearsnap@pm@chain
    \fi
  \fi}
\def\clearsnap@pm@read{\futurelet\clearsnap@pm@la\clearsnap@pm@read@}
\def\clearsnap@pm@read@{%
  \ifx\clearsnap@pm@la\clearsnap@pm@end
    \let\clearsnap@pm@nx\clearsnap@pm@atend
  \else\ifx\clearsnap@pm@la\bgroup
    \let\clearsnap@pm@nx\clearsnap@pm@fail
  \else\ifx\clearsnap@pm@la\egroup
    \let\clearsnap@pm@nx\clearsnap@pm@fail
  \else\ifx\clearsnap@pm@la\@sptoken
    \let\clearsnap@pm@nx\clearsnap@pm@fail
  \else
    \let\clearsnap@pm@nx\clearsnap@pm@take
  \fi\fi\fi\fi
  \clearsnap@pm@nx}
\def\clearsnap@pm@fail#1\clearsnap@pm@end{\let\clearsnap@pm@good\@empty}
\def\clearsnap@pm@atend\clearsnap@pm@end{\clearsnap@pm@go E\relax}
% A token's type: D digit, P period, M minus, A plus, L letter, R control sequence or active character, X other.
\def\clearsnap@pm@take#1{%
  \def\clearsnap@pm@ty{X}%
  \ifcat\relax\noexpand#1\def\clearsnap@pm@ty{R}\else
    \ifcat\noexpand#1-%
      \ifnum`#1>47 \ifnum`#1<58 \def\clearsnap@pm@ty{D}\fi\fi
      \if\noexpand#1.\def\clearsnap@pm@ty{P}\fi
      \if\noexpand#1-\def\clearsnap@pm@ty{M}\fi
      \if\noexpand#1+\def\clearsnap@pm@ty{A}\fi
    \fi
    \ifnum`#1>96 \ifnum`#1<123 \ifcat\noexpand#1a\def\clearsnap@pm@ty{L}\fi\ifcat\noexpand#1-\def\clearsnap@pm@ty{L}\fi\fi\fi
  \fi
  \clearsnap@pm@go\clearsnap@pm@ty#1}
\def\clearsnap@pm@go#1#2{%
  \ifcsname clearsnap@pm@s\clearsnap@pm@st @#1\endcsname
    \csname clearsnap@pm@s\clearsnap@pm@st @#1\endcsname#2%
  \else
    \let\clearsnap@pm@good\@empty
  \fi
  \ifx\clearsnap@pm@good\relax
    \if#1E\else\expandafter\expandafter\expandafter\clearsnap@pm@read\fi
  \else
    \if#1E\else\expandafter\expandafter\expandafter\clearsnap@pm@fail\fi
  \fi}
\def\clearsnap@pm@add#1{\expandafter\def\expandafter\clearsnap@pm@term\expandafter{\clearsnap@pm@term#1}}
\def\clearsnap@pm@set#1{\def\clearsnap@pm@st{#1}}
% Transitions: start, term start, integer digits, after a leading period, fraction digits, a unit letter, after a
% unit, after a register.
\@namedef{clearsnap@pm@sA@M}#1{\def\clearsnap@pm@sign{-}\clearsnap@pm@set B}
\@namedef{clearsnap@pm@sA@D}{\@nameuse{clearsnap@pm@sB@D}}
\@namedef{clearsnap@pm@sA@P}{\@nameuse{clearsnap@pm@sB@P}}
\@namedef{clearsnap@pm@sA@R}{\@nameuse{clearsnap@pm@sB@R}}
\@namedef{clearsnap@pm@sB@D}#1{\clearsnap@pm@add#1\clearsnap@pm@set N}
\@namedef{clearsnap@pm@sB@P}#1{\clearsnap@pm@add#1\clearsnap@pm@set Q}
\@namedef{clearsnap@pm@sB@R}#1{\clearsnap@pm@treg#1}
\@namedef{clearsnap@pm@sN@D}#1{%
  \def\clearsnap@pm@tmp{0}\ifx\clearsnap@pm@term\clearsnap@pm@tmp\let\clearsnap@pm@good\@empty\fi
  \clearsnap@pm@add#1}
\@namedef{clearsnap@pm@sN@P}#1{\clearsnap@pm@add#1\clearsnap@pm@set F}
\@namedef{clearsnap@pm@sN@L}#1{\clearsnap@pm@add#1\def\clearsnap@pm@unit{#1}\clearsnap@pm@set U}
\@namedef{clearsnap@pm@sN@R}#1{\clearsnap@pm@treg#1}
\@namedef{clearsnap@pm@sN@A}#1{\clearsnap@pm@endterm n\def\clearsnap@pm@op{}\clearsnap@pm@set B}
\@namedef{clearsnap@pm@sN@M}#1{\clearsnap@pm@endterm n\def\clearsnap@pm@op{-}\clearsnap@pm@set B}
\@namedef{clearsnap@pm@sN@E}#1{\clearsnap@pm@endterm n}
\@namedef{clearsnap@pm@sQ@D}#1{\clearsnap@pm@add#1\clearsnap@pm@set F}
\@namedef{clearsnap@pm@sF@D}#1{\clearsnap@pm@add#1}
\@namedef{clearsnap@pm@sF@L}{\@nameuse{clearsnap@pm@sN@L}}
\@namedef{clearsnap@pm@sF@R}#1{\clearsnap@pm@treg#1}
\@namedef{clearsnap@pm@sF@A}{\@nameuse{clearsnap@pm@sN@A}}
\@namedef{clearsnap@pm@sF@M}{\@nameuse{clearsnap@pm@sN@M}}
\@namedef{clearsnap@pm@sF@E}{\@nameuse{clearsnap@pm@sN@E}}
\@namedef{clearsnap@pm@sU@L}#1{%
  \ifcsname clearsnap@pm@u@\expandafter\string\clearsnap@pm@unit\string#1\endcsname\else\let\clearsnap@pm@good\@empty\fi
  \clearsnap@pm@add#1\clearsnap@pm@set V}
\@namedef{clearsnap@pm@sV@A}#1{\clearsnap@pm@endterm u\def\clearsnap@pm@op{}\clearsnap@pm@set B}
\@namedef{clearsnap@pm@sV@M}#1{\clearsnap@pm@endterm u\def\clearsnap@pm@op{-}\clearsnap@pm@set B}
\@namedef{clearsnap@pm@sV@E}#1{\clearsnap@pm@endterm u}
\@namedef{clearsnap@pm@sR@A}{\@nameuse{clearsnap@pm@sV@A}}
\@namedef{clearsnap@pm@sR@M}{\@nameuse{clearsnap@pm@sV@M}}
\@namedef{clearsnap@pm@sR@E}{\@nameuse{clearsnap@pm@sV@E}}
% A register term (after a factor or not) is a dimen register; the guard checks it is still the same one.
\def\clearsnap@pm@treg#1{%
  \edef\clearsnap@pm@tmp{\noexpand\clearsnap@pm@dimtest\meaning#1\string\dimen\noexpand\clearsnap@pm@end}\clearsnap@pm@tmp
  \ifx\clearsnap@pm@isdim\relax
    \global\expandafter\let\csname clearsnap@pm@d@\string#1\endcsname#1%
    \edef\clearsnap@pm@guard{\unexpanded\expandafter{\clearsnap@pm@guard}\noexpand\clearsnap@pm@chk\noexpand#1%
      \expandafter\noexpand\csname clearsnap@pm@d@\string#1\endcsname}%
    \clearsnap@pm@add#1\clearsnap@pm@set R
  \else
    \let\clearsnap@pm@good\@empty
  \fi}
\begingroup\edef\clearsnap@pm@tmp{\endgroup
  \def\noexpand\clearsnap@pm@dimtest##1\string\dimen##2\noexpand\clearsnap@pm@end{%
    \noexpand\clearsnap@pm@dimtest@{##1}##2x\noexpand\clearsnap@pm@end}}\clearsnap@pm@tmp
\def\clearsnap@pm@dimtest@#1#2#3\clearsnap@pm@end{%
  \let\clearsnap@pm@isdim\@empty
  \def\clearsnap@pm@tmp{#1}%
  \ifx\clearsnap@pm@tmp\@empty \ifnum`#2>47 \ifnum`#2<58 \let\clearsnap@pm@isdim\relax\fi\fi\fi}
% A term goes into the code: the first as \pgfmath@dimen=, a later one with \advance (n: a number, read as pt).
\def\clearsnap@pm@pt{pt}
\def\clearsnap@pm@endterm#1{%
  \edef\clearsnap@pm@tmp{\unexpanded\expandafter{\clearsnap@pm@term}\if#1n\clearsnap@pm@pt\fi\relax}%
  \ifx\clearsnap@pm@terms\@empty
    \edef\clearsnap@pm@code{\noexpand\pgfmath@dimen=\clearsnap@pm@sign\unexpanded\expandafter{\clearsnap@pm@tmp}}%
    \ifx\clearsnap@pm@sign\@empty\else
      \edef\clearsnap@pm@guard{\unexpanded\expandafter{\clearsnap@pm@guard}\noexpand\clearsnap@pm@chk
        \noexpand\pgfmathneg@\noexpand\clearsnap@pm@expectneg}%
    \fi
  \else
    % (pgfmath's add and subtract end by assigning the sum to itself, which keeps it within \maxdimen.)
    \edef\clearsnap@pm@code{\unexpanded\expandafter{\clearsnap@pm@code}\noexpand\advance\noexpand\pgfmath@dimen by%
      \clearsnap@pm@op\unexpanded\expandafter{\clearsnap@pm@tmp}\noexpand\pgfmath@dimen\noexpand\pgfmath@dimen}%
    \edef\clearsnap@pm@guard{\unexpanded\expandafter{\clearsnap@pm@guard}\noexpand\clearsnap@pm@chk
      \ifx\clearsnap@pm@op\@empty\noexpand\pgfmathadd@\noexpand\clearsnap@pm@expectadd
      \else\noexpand\pgfmathsubtract@\noexpand\clearsnap@pm@expectsubtract\fi}%
  \fi
  \edef\clearsnap@pm@terms{\clearsnap@pm@terms#1}%
  \let\clearsnap@pm@term\@empty \let\clearsnap@pm@unit\@empty}
% A register is one control word (the text is just its name and a space), or - and one, that is a dimen register.
\def\clearsnap@pm@one{\clearsnap@pm@stop}
\def\clearsnap@pm@rsplit#1#2\clearsnap@pm@end{\def\clearsnap@pm@first{#1}\def\clearsnap@pm@rest{#2}}
\def\clearsnap@pm@regscan{%
  \expandafter\clearsnap@pm@rsplit\pgfmath@expression\clearsnap@pm@stop\clearsnap@pm@end
  \ifx\clearsnap@pm@first\pgfmath@char@minus
    \let\clearsnap@pm@reg@kind\clearsnap@pm@negreg
    \expandafter\clearsnap@pm@rsplit\clearsnap@pm@rest\clearsnap@pm@end
  \else
    \let\clearsnap@pm@reg@kind\clearsnap@pm@reg
  \fi
  \ifx\clearsnap@pm@rest\clearsnap@pm@one
    \edef\clearsnap@pm@tmp{\ifx\clearsnap@pm@reg@kind\clearsnap@pm@negreg-\fi\detokenize\expandafter{\clearsnap@pm@first}}%
    \ifx\clearsnap@pm@tmp\clearsnap@pm@key
      \expandafter\clearsnap@pm@isdimen\clearsnap@pm@first
    \fi
  \fi}
\ExplSyntaxOn
\regex_const:Nn \c__clearsnap_pm_dimen_regex { \A \\dimen \d+ \Z }
\cs_new_protected:Npn \clearsnap@pm@isdimen #1
  {
    \exp_args:NNe \regex_match:NnT \c__clearsnap_pm_dimen_regex { \token_to_meaning:N #1 }
      {
        \cs_gset_eq:cN { clearsnap@pm@r@ \clearsnap@pm@key } #1
        \cs_set_eq:NN \clearsnap@pm@kind \clearsnap@pm@reg@kind
      }
  }
\ExplSyntaxOff
\AtBeginDocument{\clearsnap@pm@install}
\makeatother
